\chapter*{Conclusion}
\addcontentsline{toc}{chapter}{Conclusion}
\markboth{Conclusion}{Conclusion}

\section*{Récapitulatif des réalisations}

Au terme de ce stage d'un mois, l'application web de gestion des appels d'offres
a été développée, testée et déployée dans son intégralité. Les objectifs fixés en
début de stage ont été atteints :

\begin{itemize}
    \item Authentification sécurisée JWT avec cookie httpOnly scopé,
    contrôle d'accès par rôles (RBAC) sur l'ensemble des endpoints ;
    \item Cycle de vie complet des appels d'offres (DRAFT → PUBLISHED → EVALUATION
    → AWARDED) avec machine à états côté serveur ;
    \item Portail public permettant aux sociétés externes de déposer leurs propositions
    sans compte, avec rate limiting et vérification des deadlines ;
    \item Chiffrement AES-256-GCM sélectif des documents sensibles,
    déchiffrement transparent au téléchargement ;
    \item Notifications automatiques par le scheduler Docker ;
    \item Export PDF des synthèses d'offres ;
    \item Journal d'audit complet sur toutes les opérations d'écriture ;
    \item Pipeline CI GitHub Actions (lint + tests) actif sur chaque push.
\end{itemize}

\section*{Difficultés rencontrées}

\paragraph{Isolation du portail public} La principale difficulté architecturale
a été de concevoir le portail public comme une surface réellement isolée de l'API
interne. Le problème du cookie de rafraîchissement -- qui, sans restriction de chemin,
aurait été envoyé sur toutes les requêtes du portail -- a nécessité de restreindre
le cookie à \texttt{path=/api/auth/}, une subtilité non documentée dans la plupart
des tutoriels JWT.

\paragraph{Chiffrement sélectif} L'intégration du chiffrement AES-256-GCM a
soulevé la question de la gestion des clés : stocker la clé dans le code source
ou dans \texttt{docker-compose.yml} aurait constitué une faille évidente.
La solution par variable d'environnement (\texttt{DOCUMENT\_ENCRYPTION\_KEY})
et l'exclusion des fichiers \texttt{.env} du dépôt Git ont résolu ce problème.

\paragraph{Pipeline CI} Les premiers commits ont déclenché des échecs de pipeline
liés à des migrations manquantes et des violations de lint. La correction a nécessité
de comprendre le flux de migration Django dans un environnement de test isolé.

\section*{Compétences acquises}

Ce stage a été l'occasion de mettre en pratique et d'approfondir plusieurs
compétences :

\begin{itemize}
    \item Développement d'API REST sécurisées avec Django REST Framework ;
    \item Sécurité web applicative (OWASP Top 10, JWT, RBAC, chiffrement) ;
    \item Conteneurisation avec Docker Compose et architecture multi-services ;
    \item Intégration continue avec GitHub Actions ;
    \item Conception de systèmes : modèle de données, machine à états, audit.
\end{itemize}

\section*{Perspectives}

L'application est fonctionnelle et déployable, mais plusieurs améliorations
pourraient être envisagées pour une mise en production réelle :

\begin{description}
    \item[Certificat SSL valide] Remplacer le certificat auto-signé par un certificat
    Let's Encrypt pour éliminer les avertissements navigateur.
    \item[Notifications par email] Ajouter l'envoi d'emails (SMTP) en complément
    des notifications en application.
    \item[Rotation des clés AES] Implémenter un mécanisme de re-chiffrement
    des documents lors d'un changement de clé.
    \item[Optimistic locking] Prévenir les conflits d'édition simultanée
    sur les offres via un champ de version.
    \item[Docker Secrets] Remplacer les variables d'environnement par Docker Secrets
    ou un gestionnaire de secrets dédié (Vault) en production.
\end{description}

\todo[inline, color=orange!30]{À RÉDIGER : paragraphe de clôture personnel — bilan humain, expérience globale du stage, recommandation}


}
\end{spacing}

\newpage
